\section{Introduction}

A quantifier can be read as a move. In the evaluation game of first-order
logic \cite{Hintikka1973}, one player supplies witnesses for $\exists$,
another supplies them for $\forall$, and a sentence is true just when the
first player has a winning strategy. On this reading $\exists$ and $\forall$
are not really two binders. They are one binder with two possible owners.
It is then natural to ask what happens when there are more than two players.

Write $\Q_r(x:A).G$ for the game in which player $r$ chooses a value $x$ of
type $A$, in full view of everyone, after which play continues as $G$. No
player is permanently the verifier, and none is permanently the environment.
This syntax is not new: player-indexed connectives and quantifiers, together
with role permutations, appear in Abramsky's game semantics for multi-agent
logics \cite{Abramsky2007}, in the logics and algebras of Olde Loohuis and
Venema \cite{OldeLoohuisVenema2010}, and in the multirole logic of Xi and Wu
\cite{XiWu2023}. The question we take up is semantic. What does such a game
\emph{mean}, and when should two of them count as the same?

With two players and opposed objectives the answer is classical: ask who has
a winning strategy. With three players that answer says very little.
Consider the game in which player $c$ chooses whether $a$ or $b$ wins:
\[
 \at a\ch c\at b.
\]
Nobody can force a win here, and a semantics of guarantees sees nothing in
it. Yet something is plainly going on. Player $c$ loses either way, so
nothing in $c$'s interests settles the matter, and either of the others may
end up winning. Placed inside a larger game, that possibility matters.

We therefore observe rational play rather than guarantees. For a finite
public game $G$, let $B(G)$ be the set of players who win in some pure
subgame-perfect equilibrium of $G$, so that $B(\at a\ch c\at b)=\{a,b\}$.
Backward induction then gives a very simple rule. At a node owned by $r$,
with branches of values $S$ and $T$,
\[
 S\star_r T=\begin{cases}
 \{r\},&S=\{r\}\text{ or }T=\{r\},\\
 S\cup T,&\text{otherwise}.
 \end{cases}
\]
If a branch guarantees $r$'s win, $r$ takes it and nothing else survives.
Otherwise every possibility is retained. The one point that deserves a
moment's thought is that a branch in which $r$ merely \emph{might} win does
not absorb. ``Might'' refers to which equilibrium is played inside that
branch, and against an equilibrium in which $r$ loses there, $r$ is
indifferent and may go either way.

Every finite public game thus has a value among the nonempty subsets of the
players, and these values compose. This paper is a study of the resulting
algebra. It has $2^K-1$ elements and is defined in one line. For two players its
ground terms are Boolean logic in disguise. From three players on it is something else:
distributivity fails, and no representation by Boolean meets and joins
exists. Yet the algebra is far from shapeless, and its structure can be
determined rather exactly.

\paragraph{What we find.}
Our results fall into five groups.

\emph{Equations have small counterexamples.} Intermediate results of a term
may range over all winner sets, but its inputs need not. Every failed
equation already fails when each variable has one or two possible winners,
and every such input is needed. The resulting test domain, of size
$K(K+1)/2$, is optimal uniformly in the number of players and variables
(\cref{sec:pair-tests}).

\emph{Idempotency picks out a core.} A term of the full language, with winner
constants and internal role permutations, is idempotent exactly when it can
be written with owned choices and variables alone (\cref{sec:intrinsic}).
On this core each player's view has only $K+1$ states
(\cref{sec:perspectives}). A common-source theorem shows what determinism
costs: an idempotent rule that always returns a unique winner on unique-winner
inputs can do nothing there except project one of them.

\emph{One small fragment is completely understood.} The binary operations of
that kind are the projections and \emph{reset}: keep a resolved right operand,
otherwise take union (\cref{sec:reset}). Reset has polynomial equality and a
representation by antimatroids. Its three-element generator turns out to be a
known order algebra, whose published finite basis gives a complete calculus.

\emph{Tractability ends abruptly, and at an exact place.} Equality is
coNP-complete. Owned choice, union, and reset are each easy alone, and every
pair of them is hard (\cref{sec:complexity}). Yet reset with one owner's
guarded choice stays polynomial and has a minimum generator of four elements.
For an individual observer, hardness begins exactly when two \emph{other}
owners are present (\cref{sec:one-owner}).

\emph{Beyond public choice, probabilities cannot be avoided.} Commit--reveal
rounds with complete disclosure compose exactly, but on sets of winner
distributions rather than winner sets. With three players, no finite
winner-set algebra can reach them (\cref{sec:sealed-rounds}).

The appendices also contain classifications that the main text only touches on: all
congruences, which internal switching reduces to a chain of ambiguity thresholds; the $751$
three-role unary operations; and the general ambiguity-action quotient with its two-permutation
fibers.

\paragraph{What we do not claim.}
These are results about one specified semantics. Apart from the calculus for
reset, which is imported, we give no proof system, and a finite equational
basis for the full algebra or its core remains open. Nor do we give a
dependent type theory, or a cut theorem for hidden games. The round recursion
requires full public disclosure at its boundaries.
\Cref{app:dependent} sets out interfaces that such a theory would have to
respect: reducing full strategies to executable plans preserves execution
but can lose subgame perfection, and publishing dependent data requires
stable binding and a permitted local context.

\paragraph{Relation to existing work.}
Olde Loohuis and Venema's algebras observe winning strategies; our observation
is possible SPE winners with all ties retained, and individual forcing does
not define it. Xi and Wu's multirole logic has multiparty cut elimination and
is the closest proof-theoretic precursor to our binder. Full contingent
strategies and dependent game trees have established treatments
\cite{Rubinstein1991,EscardoOliva2023}, and finite backward induction supplies
the starting algebra. Multiparty coherence \cite{CarboneEtAl2016} and dependent
command/response interaction \cite{HancockSetzer2000} are further foundations
for the dependent interface.

We rely on established mathematics throughout, and we use it directly: the
classification $\operatorname{Pol}(\iota_3)=S_2$ for atomic expressivity
\cite{Szendrei2012}; the factorpower of the symmetric group for ambiguity
actions \cite{GanyushkinMazorchuk1995,Mazorchuk2009}; standard contextual
congruences and closure operators \cite{BlokPigozzi1989,Font1993,Crapo1967};
the order algebra, finite basis, related path normal forms, and
residual-largeness result of Freese et al.\ \cite{FreeseEtAl2002}; the known
minimal clone of Waldhauser \cite{Waldhauser2006}; and the dependency
structures with choice of Bazerman et al.\ \cite{BazermanMinichielloPuzio2023}.
None of these is counted here as a new classification.

The contribution is the structure of possible-SPE-winner semantics: sharp
tests, intrinsic syntactic elimination, exact tractable and hard fragment
boundaries, and the transition from finite to infinite at sealed rounds. The
prerequisite tables specialize the older path forms to mixed winner-set
valuations, with explicit algorithms and substitution
(\cref{reset:rem:priority}). The comparisons above identify what the cited
sources supply. Priority for the remaining guard and forest connections is not
established by those comparisons.

Finally, the probabilistic open-game category proposed by Ghani et al.\
\cite{GhaniEtAl2020} has finite counterexamples to its unit and associativity
laws \cite{GameTheoryOpenGameObstructions2026}, which we reproduce in
\cref{se:sec:open-game-obstruction}. Our round recursion is proved directly and
does not use that composition.

\paragraph{How to read this paper.}
The main text states the results and explains why they hold, giving complete
arguments where these are short. Longer proofs are collected in
\cref{app:proofs}. The remaining appendices contain the full unary, congruence,
ambiguity-monoid, and forest-model developments, the open-game obstruction, the
dependent interface, and the details of our finite checks. \Cref{tab:status}
records, for each main result, whether it is kernel-checked in Lean, proved on
paper, or imported from the literature.
